\section{An obstruction for actual counting trajectories}
Structural preservation is stronger than uniformity of one actual trajectory. The following result addresses the latter without requiring positive counts in every residue.
\begin{theorem}[Finite-range mass obstruction]\label{thm:mass}
Let $w:\Z\to\mathbb R_{\ge0}$ have finite nonempty support $S$, with $m=\min S$, $M=\max S$. Let $I_n=[A_n,B_n]\cap\Z$ be finite integer intervals, and put
\[
J_n=I_n\cap[A_{n-1}+M,B_{n-1}+m].
\]
Suppose actual rows $p_b,\ldots,p_h$ are $q$-residue-uniform, every interval in this block has at least $q$ vertices, and each $J_n$, $b<n\le h$, has at least $q$ consecutive vertices. If
\[
\left\lfloor |I_h|/q\right\rfloor>\left\lceil |I_b|/q\right\rceil,
\]
then $p_b=0$.
\end{theorem}
\begin{proof}
Let $c_n(s)$ be the common value in residue $s$, $C_n=\sum_s c_n(s)$, and $P_n=\sum_{y\in I_n}p_n(y)$. For each residue choose a target in $J_n$. All supported predecessors are admitted, so
\[
c_n(s)=\sum_d w(d)c_{n-1}(s-d),\qquad C_n=K C_{n-1},\quad K=\sum_d w(d)>0.
\]
Each word has total outgoing extension weight at most $K$, hence $P_n\le K P_{n-1}$. If $C_b>0$, residue multiplicities give
\[
\lfloor |I_h|/q\rfloor C_h\le P_h\le K^{h-b}P_b
\le K^{h-b}\lceil |I_b|/q\rceil C_b.
\]
Substitute $C_h=K^{h-b}C_b$ and cancel, contradicting the strict inequality. If $C_b=0$, nonnegativity gives $p_b=0$.
\end{proof}

\begin{proposition}\label{prop:cones}
If $|I_n|\to\infty$ and the full-predecessor cores $J_n$ eventually contain $q$ consecutive vertices, eventual residue uniformity forces extinction. In particular this holds when interval widths diverge and successive endpoint movements are bounded, including rational cones with bounded periodic boundary offsets.
\end{proposition}
\begin{proof}
Choose any sufficiently late $b$ and then $h$ giving the strict multiplicity gap; apply Theorem~\ref{thm:mass}. For bounded endpoint movements $D_A,D_B$, the size of $J_n$ is at least $|I_{n-1}|-C$, where $C=\max(D_A,M)-\min(-D_B,m)$. Diverging widths give the core condition. Rounded affine boundaries of distinct slopes have diverging widths and bounded endpoint movements, as do their bounded periodic perturbations.
\end{proof}

\subsection{A shorter triangular certificate}
\begin{theorem}[Boundary propagation]\label{thm:short}
In the triangle, let $L=\max(0,M)$, $R=\max(0,1-m)$ and $v=\max(M,1-m)\ge1$. If
\[
b\ge\max(q-1,L+R+q-2),\qquad T=\lceil q/v\rceil,
\]
then $q$-uniform rows $b,\ldots,b+T$ force $p_b=0$. This holds also for arbitrary prescribed nonnegative data at layer $b$.
\end{theorem}
\begin{proof}
The full cores give the residue convolution in Theorem~\ref{thm:mass}. At any other target, the actual sum omits terms from that full convolution. Nonnegativity and equality of the two values force every omitted term to vanish.

Take an active source residue and a supported step $d>0$. Its least representative $r\in[0,q-1]$ must satisfy $r+d<q$: otherwise the least target of residue $r+d$ has a negative predecessor and a positive omitted term. Full convolution propagates positivity, so repeated use of $d$ increases the representative by $d$ without wrapping. This is impossible after $\lceil q/d\rceil$ transitions. For $d\le0$ use the least distance $\delta$ from the moving top to the active residue. It increases by $1-d$ without wrapping, since a wrap would give a predecessor above the previous top. This is impossible after $\lceil q/(1-d)\rceil$ transitions. Choose $d=M$ or $m$ to attain $v$.
\end{proof}
For restarted data the bound is sharp for every $q,v$: take the single jump $d=1-v$, a sufficiently large $b$ divisible by $q$, and initial value one exactly on residue zero. Translation gives precisely $\lceil q/v\rceil$ consecutive nonzero uniform rows before a missing top representative breaks uniformity. This is not an origin-start optimality example. In the natural origin-start class $\operatorname{supp}f\subseteq\{0,1\}$, a sharper fact holds: no row $n\ge q+1$ is $q$-uniform. If both weights $a,b>0$, comparing binomial counts at $0,q$ and $1,q+1$ would require
$\binom n{q+1}/\binom nq=n$, whereas it equals $(n-q)/(q+1)$. Single-step cases fail already for $n\ge q$; equal weights attain uniformity at $n=q$.

The certificate also excludes proposed inverse models: $t+1$ nonzero uniform rows satisfying their core threshold imply $tv<q$. In particular two such rows rule out every finite-support candidate with $v\ge q$, without recovering its weights. The threshold depends on the proposed support; this is not an unconditional bound on an unknown infinite kernel.

\paragraph{Assumption boundaries.}
For layer-dependent nonnegative kernels supported in one common finite interval, Theorem~\ref{thm:mass} still holds: replace $K^{h-b}$ by $\prod_{n=b+1}^h K_n$. Zero $K_n$ causes extinction. Thus eventual obstruction does not need time homogeneity. Without a common range bound, $f_n=1_{[1-n,n]}$ produces $p_n(y)=n!$. The same counts arise for the infinite homogeneous kernel $f\equiv1$. Bounded layers $I_n=\{0\}$ with $f(0)=1$ also persist. Signed kernels and sparse nonconsecutive uniform rows remain open extensions here; failure of our proof is not a counterexample. Standard killed-walk operators on fixed sets \cite[Section 12.4.1]{LawlerLimic} provide related background, not a priority determination for this expanding-window certificate.
